\subsection{Opening work carries both cost and completed effort}
\label{sec:opening-production-cohorts}

The earlier production example began with an empty factory account. An opening
work-in-progress balance changes both the cost available and the work still
needed. Materials may already be present even though conversion remains
unfinished. Treating the opening units as entirely new would charge that
effort twice; ignoring their costs would lose part of the inventory.

Consider homogeneous units with all materials added at the start and
conversion effort incurred evenly. Ten opening units contain material
cost 100 and conversion cost 50 and are halfway through conversion.
The factory starts another 90 units, incurring material cost 900 and
conversion cost 510. It completes 80 units and leaves 20 units halfway
through conversion. Assume the opening units are completed first, no
spoilage, no idle-capacity cost, and net realizable values above cost.
Physical units reconcile as \(10+90=80+20\); available cost is
\(100+50+900+510=1560\). These are supplied production facts, not estimates
derived from a reported inventory total.

Under FIFO process costing, retain the opening 150 and allocate current
cost only to current work. Completing the ten opening units requires
five equivalent conversion units. Seventy units are started and completed,
and the closing twenty require ten equivalent conversion units. Current
conversion effort is therefore \(5+70+10=85\), at \(510/85=6\) per
equivalent unit. Current material cost is \(900/90=10\) per new unit.
Completing opening work costs \(150+5\times6=180\); the seventy new
completed units cost \(70\times(10+6)=1120\). Transfer 1300 from WIP
to finished goods and retain \(20\times10+10\times6=260\) in WIP.
Completion is not a sale, so this transfer does not itself create cost
of sales or a cash payment.

Weighted-average process costing pools the opening and current costs.
Material cost remains 10 per equivalent unit, but conversion cost becomes
\((50+510)/(80+10)=56/9\). Completed goods then cost \(11680/9\)
and closing WIP \(2360/9\), still totaling 1560. Posting at cents gives
1297.78 and 262.22 with a preserved rounding allocation. The methods
distribute the same total cost differently because the opening conversion
cost differs from the current-period rate. Neither method is selected
merely to obtain a preferred closing inventory. IAS~2.23--27 requires
specific identification where applicable and consistent FIFO or
weighted-average treatment for the relevant interchangeable inventories;
the equivalent-unit calculations are the local implementation of those
cost policies \citep[IAS~2.10--16, 23--27, 34]{eu20231803}.

For several stages, transfer the preceding stage's reconciled cost into
the next stage as a separate cost component. Record when that component
enters the process; it need not follow the completion percentage of later
materials or conversion. Each stage retains opening units and costs,
transfers in, completed transfers out, closing units, progress by component,
inspection points, losses and recoveries. An impossible completion fraction,
unexplained physical loss, missing opening cost or inconsistent transfer
amount blocks the cost allocation rather than becoming an unexplained
unit-cost adjustment.

\subsubsection{Abnormal loss cannot become a higher cost of the good units}
\label{sec:production-spoilage}

Normal process loss and abnormal waste need different treatment. IAS~2.16(a)
excludes abnormal wasted materials, labour and other production costs from
inventory. Its normal-capacity rule also prevents idle fixed overhead from
being hidden in a higher unit cost \citep[IAS~2.12--16, 21]{eu20231803}.
The normal-loss estimate must therefore be supported before observing the
actual loss, with its inspection stage and recoveries specified.

In a separate fully processed batch, 100 input units cost 1000. Assume
normal yield is 95 units, all lost units have consumed the same full effort
at inspection, there is no salvage, and this cost pool contains no fixed
overhead. The normal unit cost is \(1000/95=200/19\). If actual good
output is 90, five units of loss are normal and five abnormal. Good
inventory receives \(90\times200/19=18000/19\), while abnormal loss
of \(1000/19\) is expensed. At cents these are 947.37 and 52.63,
which reconcile to 1000. Allocating all 1000 to the ninety good units
would conceal the abnormal loss. Earlier inspection, partial conversion,
recoverable scrap or a different supported normal yield changes the
calculation; those facts cannot be supplied by the final good-unit count.

\subsection{A pension surplus is useful only to the extent it is available}
\label{sec:pension-funding-availability}

The preceding pension examples supplied a recoverable economic benefit.
IFRIC~14 specifies how the plan's terms, statutory requirements and funding
basis constrain that benefit. Separate an unconditional right to a refund
from the ability to reduce future contributions, and separate funding for
past service from funding for future service. An IAS~19 surplus can coexist
with a shortfall measured on the funding basis
\citep[IAS~19.64; IFRIC~14.7--24]{eu20231803}.

Start with plan assets 130 and defined-benefit obligation 100, giving
surplus 30. Suppose the employer has an unconditional right to the surplus
but obtaining it costs the plan 5. The available refund is 25 and the
recognized asset is 25, with a ceiling adjustment of 5. Costs depend on
the permitted refund route, including settlement costs if winding up the
plan is necessary. A right conditional on an uncertain event outside the
employer's complete control does not establish this refund benefit.
For a refund defined as the whole surplus or a proportion of it,
IFRIC~14.15 makes no additional time-value adjustment to that refund
\citep[IFRIC~14.11--15]{eu20231803}.

Contribution reductions require a different calculation. Use employer
service costs, excluding employee-funded amounts, over the shorter of
expected plan life and entity life. Assumptions must be consistent with
the benefit obligation and the situation at the reporting date.
In a two-period example with zero discounting solely for illustration,
future employer service costs are 20 and 20. Without a prepayment,
required future-service funding would be 12 and 25. Their net difference
is \(8-5=3\). Suppose a prior prepayment separately reduces future
required contributions by 7. The available reduction benefit is then
\[
 7+\max\{0,(20-12)+(20-25)\}=10.
\]
The floor applies to the combined paragraph 20(b) amount, not separately
to each period: ignoring the second period's excess would incorrectly
produce 15. With no minimum funding requirement, use the available future
employer service costs under paragraph 16 instead
\citep[IFRIC~14.16--22]{eu20231803}.

The refund and contribution routes may be mutually exclusive. In that
case the available benefit is the greater of 25 and 10, not their sum 35;
the recognized asset is also capped by the actual surplus. If the refund
right fails but the contribution route remains supported, the asset is
10 and the ceiling adjustment is 20. Compatible combinations require
their actual feasible terms and cannot consume the same surplus twice.

Past-service funding can require a liability before any payment.
Suppose plan assets are 120 and the IAS~19 obligation is 100, but none
of the surplus can ever be refunded or reduce contributions. The
recognized pension asset is zero. A separate funding-basis assessment
now requires a past-service contribution of 10 that will also be wholly
unavailable after payment. IFRIC~14.23--24 requires the resulting
liability when that obligation arises. With no intervening interest or
other pension changes, recognize an OCI ceiling loss of 10 and a net
pension liability of 10. On payment, debit that liability and credit
cash 10. Plan assets become 130, but the recognized asset remains zero
and payment causes no second loss or gain. Starting company cash and
equity of 100 become cash 90 and equity 90. The obligation's recognition
and the payment are distinct events
\citep[IFRIC~14.23--24; IAS~19.120(c), 126--127]{eu20231803}.

\subsection{Restoration cash follows the activity being paid for}
\label{sec:restoration-cash-classification}

Return to the exhausted asset whose restoration liability is settled for
36.3. To complete its cash-flow example, specify the payment facts:
the company pays wages, materials and contractors to remove the exhausted
asset and restore the site; it acquires no new long-term asset and repays
no borrowing. The proposed classification for this conditional case is
an operating outflow of 36.3. This follows the activity definitions and
supplier/employee examples in IAS~7.6 and 14. The earlier noncash creation
of an asset and provision is not a cash outflow, and unwinding the provision
does not by itself establish that part of the eventual contractor payment
is interest paid \citep[IAS~7.6, 12--16, 31--33, 43--44;
IFRIC~1.8]{eu20231803}.

The competing interpretation classifies the payment as investing because
the provision's initial measurement entered the asset's cost. The
July 2012 IFRS Interpretations Committee staff paper examines comparable
rehabilitation facts, presents both approaches and favors operating
classification. That is support for the proposed reasoning, not a binding
Interpretation or proof of universal treatment
\citep[Appendix~A, A35--A43; status notice on p.~1]{ifric2012cashclassification}.
The design records the chosen interpretation and its facts so that a
different admitted rule can change the category without rewriting the
legal journal or total cash.

Under the proposed branch, opening cash 100 plus operating cash flow
minus 36.3, with no investing or financing flow, gives closing cash 63.7.
In year 2 the indirect reconciliation is loss 74.3, plus noncash
depreciation 71 and unwinding 3.3, less provision settlement 36.3, giving
operating outflow 36.3. Under the competing investing branch, operating
cash flow is zero and investing outflow 36.3; closing cash is identical.
The classification difference is economically visible even though both
cash bridges balance.

A new productive asset, actual borrowing repayment, separately evidenced
interest payment, or restoration caused by current inventory production
requires its own assessment. The proposed operating branch completes the
specified exhausted-asset example; it is not a mandatory default for every
restoration liability or evidence of an actual issuer's compliance.
All three additions require explicit successor records and exchanges.
They preserve the closed \texttt{CorporateJointHistory/1} profile rather
than silently enlarging its accounts or event vocabulary.
